\section{Introduction}
\label{sec:intro}

A recurring question about post-training is whether it teaches a model
anything or only selects a style it already had. The superficial alignment
hypothesis \citep{zhou2023lima} takes the second view, and output-level
evidence such as in-context alignment of base models \citep{lin2024urial}
supports it. A natural place to look for weight-level evidence is the update
itself, $\dW = W_{\text{after}} - W_{\text{before}}$: if alignment only picks
a style, perhaps $\dW$ should be small and structured, ideally low rank.

That question has to be answered with some measure of rank, and full-rank
dense matrices make the choice of measure matter. Three measures appear in
recent work: counting singular values above a fraction of the largest,
counting how many are needed to hold most of the energy
\citep{biderman2024lora}, and the entropy effective rank of
\citet{roy2007effective}, used by \citet{wu2025timber} to argue that
post-training is superficial. All three summarize the spectrum by how many
directions it spans. None of these uses compares the number against what a
matrix with no structure at all would give.

We set out to measure the rank of SFT and DPO updates on a small model and ran
into exactly this problem. By the threshold count, our SFT update was at 99.1\%
of full rank and the difference between two random initializations was at
99.6\%. Read at face value, that says the update looks like noise. Stable rank
\citep{rudelson2007sampling} on the same spectra says the opposite: the update
is three times more concentrated than the random pair, and its top direction
carries eight times the share of energy.

This paper is about that gap. Our contributions are:

\begin{itemize}
  \item \textbf{A noise baseline for every matrix.} We compare four rank
  measures on real updates against random matrices of the same shape and show
  that the threshold count, the 90\%-energy rank and the entropy effective
  rank all land within a few percent of the noise value, while stable rank
  separates real updates from noise by 1.6--9$\times$
  (\S\ref{sec:metrics}). The two readings are not in conflict: updates have a
  noise-like tail and a few dominant directions, and the measures weight these
  differently (\S\ref{sec:headtail}).
  \item \textbf{Controls that separate fine-tuning from alignment.}
  Token-matched fine-tunes on wikitext and Python code, with second seeds and
  paired layer-bootstrap intervals, are as concentrated as SFT on two model
  families (\S\ref{sec:controls}). Concentration of $\dW$ is a property of
  fine-tuning, not evidence about alignment.
  \item \textbf{Released models up to 8B.} The same pattern holds for seven
  official base/instruct pairs from 135M to 8B parameters (\S\ref{sec:released}).
  \item \textbf{A case study of Timber.} We reproduce the entropy effective
  rank measurement of \citet{wu2025timber} on two of its model pairs and show
  that its evidence for superficiality, an unchanged eRank of each weight
  matrix, is also produced by a random perturbation of the same norm. Its
  eRank cutoff for editing $\dW$ is not special: a cutoff that edits about
  40$\times$ more of the update gives the same GSM8K pass@$k$, and both effects
  are within a point of zero on the full test set (\S\ref{sec:timber}).
\end{itemize}

The training runs use a single 6\,GB GPU, so our own fine-tunes are small
(0.36B and 0.5B parameters, about two million tokens each). The measurement
results do not depend on training, and we check them on released models up to
8B. We report behavioural results from the same runs (LoRA rank sweep, DPO,
evals) in \S\ref{sec:behavior}, mostly as context.
